\label{sec:stand_der_forschung}

\subsection{Verfügbare Datenbasis und Korpora (Themenblock 1)}
\label{sec:datenbasis_korpora}

Einen guten Überblick über den aktuellen Stand und die verfügbaren Ressourcen für Leichte und Einfache Sprache in Deutschland geben \citet{madina2023easy}. Sie zeigen dabei vor allem auf, wie unterschiedlich die bestehenden Datensätze strukturiert sind und dass es an standardisierten Werkzeugen zur automatischen Bewertung der Texte fehlt.

Die Erforschung automatischer Textvereinfachung (Automatic Text Simplification, ATS) für das Deutsche litt lange Zeit unter einem deutlichen Mangel an Trainingsdaten. Erste Versuche, ein paralleles Korpus aufzubauen, gehen auf \citet{klaper2013klaper} mit dem \textit{KLAPER}-Korpus zurück, das jedoch mit ca. 250 parallelen Textpaaren (rund 7.000 Sätze und 70.000 Tokens) für moderne Deep-Learning-Ansätze deutlich zu klein war. 

Einen wesentlichen Fortschritt hinsichtlich des Korpusumfangs erzielten \citet{battisti2020corpus} mit einem multimodalen Korpus aus 6.217 Dokumenten (rund 211.000 Sätze, 5,5 Millionen Tokens). Trotz des Gesamtumfangs weist der Datensatz eine deutliche Asymmetrie auf: Während 5.461 Dokumente rein monolingual vorliegen, umfasst der parallel ausgerichtete Teil lediglich 378 Dokumentenpaare. Dies unterstreicht die Knappheit paralleler Trainingsdaten für überwachte Modelle im Deutschen.

Für den Bereich der Einfachen Sprache (Plain German) stellten \citet{stodden2023deplain} mit \textit{DEplain} ein paralleles Korpus vor. Dieses gliedert sich in einen Nachrichtenbereich (DEplain-APA) mit rund 500 Dokumentenpaaren (ca.\,13.000 manuell ausgerichtete Satzpaare) und einen Webbereich mit ca.\,150 manuell ausgerichteten Dokumenten (rund 2.000 Satzpaare). Ein zugehöriger Web-Harvester ermöglicht die automatische Vergrößerung des Webbereichs auf rund 750 Dokumentenpaare (ca.\,3.500 automatisch ausgerichtete Satzpaare). Die Autoren zeigten, dass auf Dokumentenebene trainierte Modelle feinkörnige semantische Übergänge teils unzureichend abbilden, weshalb präzise Satz-Alignments eine wesentliche Voraussetzung für Seq2Seq-Modelle darstellen. Allerdings unterscheidet sich das DEplain-Korpus durch seinen Fokus auf Einfache Sprache strukturell und linguistisch von den strengeren Vorgaben der Leichten Sprache.

Einen weiteren Beitrag leisteten \citet{toborek2023new} mit dem \textit{New Aligned Simple German Corpus}. Der Datensatz umfasst rund 700 Dokumentenpaare und über 10.000 teils manuell, teils automatisch berechnete Satzpaare. Das Korpus kombiniert Web-Scraping mit automatischer Satz-Ausrichtung. Allerdings fasst die Arbeit Texte der Leichten und der Einfachen Sprache unter einem gemeinsamen Begriff zusammen, was die gezielte Modellierung der spezifischen Richtlinien Leichter Sprache erschwert. Dieser Kompromiss verdeutlicht die methodische Herausforderung: Ein Korpusaufbau, der sich strikt auf Leichte Sprache beschränkt, ist durch die begrenzte Verfügbarkeit frei zugänglicher Parallelquellen im deutschsprachigen Web erschwert.

\subsection{Evaluierungsmetriken für die Vereinfachung (Themenblock 2)}
\label{sec:evaluierungsmetriken}

Die automatische Evaluierung von Textvereinfachungssystemen stellt eine methodische Herausforderung dar. Traditionell stützt sich die Lesbarkeitsbewertung im deutschen Sprachraum auf deterministische Formeln wie den \textit{Flesch Reading Ease} \citep{amstad1978verstaendlich}, die \textit{Wiener Sachtextformel} \citep{bamberger1984lesen} oder den \textit{LIX-Index} \citep{bjornsson1968readability} (Definitionen in Abschnitt~\ref{sec:linguistische_masse_lesbarkeit}). Diese Metriken erfassen die Oberflächenkomplexität anhand von Silben-, Wort- und Satzlängen. Während sie für allgemeine Sachtexte eine recheneffiziente Orientierung bieten, sind sie konzeptionell nicht darauf ausgelegt, spezifische Regeln der Leichten Sprache — wie Passivvermeidung, kurze Hauptsatzstrukturen oder die Bindestrich-Gliederung von Komposita — zu erfassen.

Klassische referenzbasierte N-Gramm-Metriken der maschinellen Übersetzung wie BLEU \citep{papineni2002bleu} und ROUGE \citep{lin2004rouge} messen die lexikalische Übereinstimmung mit einer Referenzübersetzung. Für die automatische Textvereinfachung erweisen sich diese jedoch oft als ungeeignet: Wie \citet{alva2021unsuitability} empirisch belegen, korreliert BLEU im Vereinfachungskontext nur schwach oder sogar negativ mit dem menschlichen Urteil über Einfachheit und Textqualität. Der Grund liegt darin, dass BLEU unzulässige 1:1-Kopien des komplexen Originals mit hohen Punktwerten belohnt, während legitime lexikalische Variationen und strukturelle Modifikationen (wie Satzteilungen und Umschreibungen) fälschlicherweise als Fehler abgestraft werden.

Als Alternative für die Textvereinfachung etablierte sich die Metrik \textit{SARI} (\textit{System Accessibility and Readability Index}) nach \citet{xu2016optimizing}. SARI vergleicht die Systemgenerierung $y$ simultan mit dem komplexen Ausgangstext $x$ und den Referenzübersetzungen $R$:
\begin{equation}
    \text{SARI} = \frac{F_{\text{add}} + F_{\text{keep}} + P_{\text{del}}}{3}
\end{equation}
Hierbei messen $F_{\text{add}}$ und $F_{\text{keep}}$ den $F_1$-Score für erfolgreich eingefügte bzw. beibehaltene $N$-Gramme, während $P_{\text{del}}$ die Präzision für getilgte $N$-Gramme des komplexen Originals bewertet. Dadurch honoriert SARI explizit Modifikationen, die über ein reines Kopieren des Ausgangstexts hinausgehen.

Im Laufe der Entstehungszeit dieser Arbeit wurde mit \textit{DETECT} von \citet{korobeynikova2024detect} eine neuartige, erlernbare Evaluierungsmetrik speziell für die deutsche Textvereinfachung vorgestellt. DETECT adaptiert das LENS-Framework für den deutschsprachigen Raum und zielt darauf ab, die Qualität automatischer Vereinfachungen über drei Dimensionen abzubilden: Verständlichkeit (\textit{Simplicity / Ease}), semantische Inhaltsbewahrung (\textit{Meaning Preservation}) sowie grammatikalische Korrektheit und Sprachfluss (\textit{Fluency / Textual Clarity}). Um den Mangel an manuell annotierten Trainingsdaten zu überwinden, wird das Regressionsmodell auf synthetisch durch Large Language Models generierten Qualitätsurteilen trainiert. In empirischen Validierungen auf einem eigens erhobenen Testdatensatz (\textit{SimpEvalDE}) wiesen die Autoren nach, dass DETECT stärkere Korrelationen mit menschlichen Qualitätsbewertungen erzielt als traditionelle Maße wie BLEU, SARI oder BERTScore.

Gleichzeitig verdeutlicht dieser Ansatz verbleibende Forschungslücken: Da der zugrundeliegende Evaluierungsdatensatz (\textit{SimpEvalDE}) primär auf Korpora der Einfachen Sprache (DEplain-APA und LHA-APA) aufbaut, sind spezifische Richtlinien der Leichten Sprache — wie das Vermeiden von Passiv- und Genitivkonstruktionen oder die typografische Gliederung von Komposita — nicht explizit modelliert. Zudem sind rein referenzbasierte Satzmetriken nur bedingt geeignet, um die Komplexität ganzer Dokumente referenzfrei zu bewerten oder als kontinuierliches Steuerungssignal während des Modelltrainings eingesetzt zu werden.

Obwohl das Thema Leichte Sprache in der deutschsprachigen NLP-Forschung über viele Jahre hinweg nur wenig Beachtung gefunden hat, zeigt das Erscheinen neuerer Arbeiten die zunehmende Dynamik und Relevanz dieses Forschungsfelds.

\subsection{Maschinelle Übersetzung und Optimierungsverfahren (Themenblock 3)}
\label{sec:maschinelle_textvereinfachung_alignment}
\label{sec:maschinelle_uebersetzung_modellierung}
\label{sec:dpo_textgenerierung}

Die maschinelle Übersetzung von Alltagssprache (AS) in Leichte Sprache (LS) lässt sich formal als monolinguale maschinelle Übersetzung (\textit{Monolingual Machine Translation}) modellieren. Im Gegensatz zur klassischen bilingualen Übersetzung erfordert dieser Transformationsprozess asymmetrische Strukturänderungen: Komplexe Satzgefüge müssen aufgeteilt (\textit{Sentence Splitting}), Passiv- in Aktivkonstruktionen überführt, irrelevante Textpassagen gekürzt (\textit{Compression}) und Fachbegriffe verständlich umschrieben werden.

Erste systematische Benchmarks neuronaler Übersetzungsmodelle für deutsche Leichte Sprache von \citet{saeuberli2020benchmarking} zeigten, dass Sequence-to-Sequence-Modelle grundlegende Satzkürzungen und Wortsubstitutionen erlernen können, bei komplexen syntaktischen Umstellungen jedoch stark durch die geringe Menge paralleler Trainingsdaten limitiert werden.

Bisherige Ansätze konzentrieren sich weitgehend auf isolierte Satzebenen oder stützen sich auf manuell annotierte Präferenzdaten. Eine zentrale Forschungslücke besteht darin, wie sich generative Übersetzungsmodelle auf Dokumenten- und Artikelebene trainieren und durch eine automatisierte Reward-Steuerung über gelernte Komplexitätsmetriken optimieren lassen.

